\section{Evaluation}
\label{sec:eval}

We evaluate \tool to answer four research questions:
\begin{itemize}
  \item[\rqref{1}] \textbf{Throughput \& Coverage:} Can \tool achieve high execution throughput and discover deep branch coverage across diverse software domains?
  \item[\rqref{2}] \textbf{Bandit Scheduling:} Does the UCB1 adaptive mutator scheduler discover more branch edges than unguided baseline mutation?
  \item[\rqref{3}] \textbf{Defect Discovery:} What concrete vulnerabilities does \tool uncover in production codebases?
  \item[\rqref{4}] \textbf{Overhead \& Cost:} What runtime overhead do the headless shims and linker-wrapped coverage interception introduce?
\end{itemize}

\begin{table}[t]
  \caption{Evaluated subject systems across target domains, test harnesses, cumulative campaign executions, and peak branch edge coverage.}
  \label{tab:subjects}
  \footnotesize
  \begin{tabular}{llrrlrr}
    \toprule
    Target & Domain & Version / Commit & LoC & Harnesses & Executions & Peak Cov \\
    \midrule
    \fastgltf & 3D Asset Parser & v0.9.1 (\texttt{705564d}) & 14{,}280 & 3 & 3{,}200{,}000 & 4{,}218 \\
    \raylib   & Game / Multimedia Engine & master (\texttt{e1a4079}) & 118{,}940 & 37 & 54{,}320{,}128 & 47{,}121 \\
    \drogon   & Asynchronous Web Server & v1.9.8 (\texttt{df418bf}) & 68{,}410 & 19 & 750{,}000{,}000 & 3{,}178 \\
    \midrule
    Total     & --- & --- & 201{,}630 & 59 & 807{,}520{,}128 & --- \\
    \bottomrule
  \end{tabular}
\end{table}



\subsection{Subjects}
\label{sec:subjects}
We evaluate \tool on three open-source C and C++ projects summarized in Table~\ref{tab:subjects}.
First, \textbf{\fastgltf} (14{,}280 LoC) parses glTF 2.0 3D assets.
Second, \textbf{\raylib} (118{,}940 LoC) provides a game engine with graphics and audio subsystems.
Third, \textbf{\drogon} (68{,}410 LoC) is an asynchronous web framework handling HTTP and WebSockets.
Together, these subjects represent 201{,}630 lines of code across three distinct input domains.

\subsection{Baselines}
\label{sec:baselines}
We compare \tool against three baselines.
The \emph{Flat Byte Mutator} applies unguided random bit flips without structural grammar knowledge.
The \emph{Static Uniform Scheduler} executes structural and binary operators with equal, fixed probabilities without bandit feedback.
The \emph{Socket-Based Baseline} tests \drogon over local TCP loopback sockets rather than in-memory mocks.

\subsection{Measurements}
For \rqref{1}, we measure execution throughput (iterations per second) and total branch edge coverage.
For \rqref{2}, we record branch edges discovered over time and empirical pull counts across bandit arms.
For \rqref{3}, we record verified memory errors, undefined behaviors, and upstream maintainer responses.
For \rqref{4}, we evaluate wall-clock overhead from linker wrapping and memory consumption of in-memory shims.

\subsection{Setup}
All experiments ran on an x86\_64 workstation with an AMD processor (16 logical cores) running Ubuntu 24.04 LTS under WSL2 with 32~GB of RAM.
Binaries were compiled using Clang 18.1.3 with \texttt{-fsanitize=address,undefined} and \texttt{-O2} optimization.
Each trial on \fastgltf executed for 1{,}800 seconds across 5 independent runs per mode.
Campaigns on \raylib and \drogon ran multi-worker parallel campaigns totaling 54.3 million and 750 million executions, respectively.

\subsection{\rqref{1}: Throughput and Branch Coverage}
\label{sec:rq1}

Table~\ref{tab:subjects} presents the aggregate campaign metrics.
Across all 59 target binaries, \tool executed 807{,}520{,}128 inputs.
On \fastgltf, \tool achieved a peak coverage of 4{,}218 branch edges with a median speed of 145 executions per second in hybrid mode.
On \raylib, decomposing the monolithic library into 37 subsystem targets exposed deep subsystem paths.
The fuzzer discovered 47{,}121 edges in the Image VM (\texttt{fuzz\_image\_ops}) and 39{,}057 edges in the software rasterizer (\texttt{fuzz\_render\_vm}).
It also found 22{,}177 edges in the glTF mesh loader (\texttt{fuzz\_model\_gltf}), reaching 122{,}885 edges across these master targets.
On \drogon, the in-memory mock shim (\texttt{MockTcpConnection}) bypassed operating system socket overhead, achieving up to 53{,}732 executions per second on the cookie parser and 35{,}410 executions per second on the WebSocket frame parser.
Peak coverage reached 3{,}178 edges in the end-to-end integration target (\texttt{fuzz\_kitchen\_sink}).

\begin{observation}
\label{obs:throughput}
Decoupling subsystems into headless in-memory targets enabled \tool to execute over 807 million cumulative inputs, reaching 53{,}732 executions per second on network parsers without socket allocation bottlenecks.
\end{observation}

\subsection{\rqref{2}: Adaptive Bandit Scheduling}
\label{sec:rq2}

To evaluate whether the UCB1 multi-armed bandit scheduler discovers more edges than unguided mutation, we monitored operator arm selections during live testing.
On \drogon's \texttt{fuzz\_http\_request}, starting from a cold seed corpus, the bandit evaluated four arms: \texttt{bin.libfuzzer}, \texttt{bin.intedge}, \texttt{bin.blockdup}, and \texttt{bin.fixup}.
Over 200 initial executions, the bandit directed 65 pulls to \texttt{bin.libfuzzer} (mean reward 0.1574) and 50 pulls to \texttt{bin.blockdup} (mean reward 0.0952), discovering 1{,}059 new edges.
The integer boundary arm (\texttt{bin.intedge}) received 43 pulls (mean reward 0.0674) and consistently triggered boundary conditions in \texttt{Content-Length} parsing.
On \raylib's \texttt{fuzz\_audio\_wav}, the bandit discovered 482 new edges, allocating 58 pulls to \texttt{bin.libfuzzer} (mean reward 0.0625) and 48 pulls to \texttt{bin.intedge} (mean reward 0.0340).

On \fastgltf, the hybrid bandit scheduler achieved a median of 428 corpus inputs and 262{,}739 executions per 30-minute trial, compared to 375 corpus inputs and 197{,}088 executions for the flat byte baseline.

\begin{observation}
\label{obs:bandit}
The UCB1 bandit scheduler dynamically allocated mutation pulls toward high-yield operators.
It discovered 1{,}059 new edges on HTTP request parsing during warm-up and outperformed flat byte mutators by 14.1\% in corpus expansion on 3D asset parsing.
\end{observation}

\begin{table}[t]
  \caption{Verified upstream security defects discovered by \tool across the three subject systems and their upstream resolution status.}
  \label{tab:findings}
  \resizebox{\linewidth}{!}{%
  \begin{tabular}{@{}lllllll@{}}
    \toprule
    ID & Target & Classification & CWE & Component & Trigger Condition & Upstream Status \\
    \midrule
    \textbf{F1} & \fastgltf & Out-of-Bounds Read & CWE-125 & \texttt{generateMeshIndices} & Accessor index $\ge$ count & Closed; maintainer fixed via commit~\cite{fastgltf_issue160} \\
    \textbf{R1} & \raylib & Integer Underflow & CWE-191 & \texttt{swPopMatrix()} in \texttt{rlsw.h} & \texttt{stackCounter == 0} & Acknowledged; PR \#6362 merged~\cite{raylib_issue6361} \\
    \textbf{R2} & \raylib & Out-of-Bounds Read & CWE-125 & \texttt{ExportAutomationEventList} & \texttt{event.type $\ge$ 24} & Acknowledged; maintainer fixed~\cite{raylib_issue6378} \\
    \textbf{R3} & \raylib & NULL Pointer Dereference & CWE-476 & \texttt{LoadWaveSamples()} & Large sample allocation & Acknowledged; maintainer fixed (commit 7639166)~\cite{raylib_issue6373} \\
    \textbf{D1} & \drogon & Process Abort in \texttt{noexcept} & CWE-755 & \texttt{HttpFileImpl::save} & Filename $>255$ bytes & Closed; PR \#2662 merged~\cite{drogon_pr2662} \\
    \textbf{D2} & \drogon & Out-of-Bounds Iterator & CWE-125 & \texttt{HttpFileImpl::save} & Traversal prefix \texttt{../../c} & Closed; maintainer fixed via commit~\cite{drogon_pr2656} \\
    \bottomrule
  \end{tabular}%
  }
\end{table}


\subsection{\rqref{3}: Defect Discovery}
\label{sec:rq3}

Table~\ref{tab:findings} summarizes the 6 upstream-resolved defects discovered by \tool across the three target codebases.

\paragraph{Defect F1 (\fastgltf): Heap Buffer Overflow.}
In \texttt{fastgltf::Parser::generateMeshIndices}, when processing glTF meshes lacking index buffers, the library generates sequential indices based on the vertex position accessor.
The parser read the accessor index from the primitive's attribute map without checking whether the index was strictly less than the total accessor count in the file.
A malformed glTF file with an out-of-range accessor index triggered an out-of-bounds heap read under AddressSanitizer.
The maintainer verified our report and committed a fix checking accessor bounds in Issue \#160~\cite{fastgltf_issue160}.

\paragraph{Defect R1 (\raylib): Matrix Stack Underflow.}
In \raylib's software rasterizer (\texttt{rlsw.h} and \texttt{rlgl.h}), transform matrices are managed via an internal matrix stack with an unsigned counter \texttt{stackCounter}.
Calling \texttt{swPopMatrix()} or \texttt{rlPopMatrix()} when \texttt{stackCounter == 0} decremented the counter, causing an unsigned underflow to \texttt{UINT\_MAX}.
Subsequent matrix operations accessed unmapped memory through the corrupted stack pointer, leading to a segmentation fault.
The maintainer acknowledged the bug in Issue \#6361 and merged Pull Request \#6362~\cite{raylib_issue6361}.

\paragraph{Defect R2 (\raylib): Out-of-Bounds Read in Event Export.}
In \raylib's serialization module (\texttt{src/rcore.c}), \texttt{ExportAutomationEventList()} serializes automation events to text.
The function indexes a global lookup table:
\begin{lstlisting}[numbers=none,escapechar=@]
autoEventTypeName[list.events[i].type]
\end{lstlisting}
The table contains 24 event name strings ($[0..23]$).
When reading malformed \texttt{.rae} files containing \texttt{event.type} values $\ge 24$, the function read past the bounds of the global array.
The maintainer acknowledged the defect in Issue \#6378 and committed a fix~\cite{raylib_issue6378}.

\paragraph{Defect R3 (\raylib): NULL Pointer Dereference in Audio Sample Allocation.}
In \raylib's audio subsystem (\texttt{src/raudio.c}), \texttt{LoadWaveSamples()} loads raw audio waveforms.
When parsing an audio file with an exorbitant sample count, the memory allocation call \texttt{RL\_MALLOC()} failed and returned \texttt{NULL}.
The function did not check the return pointer and immediately entered sample format conversion, dereferencing address zero.
The maintainer acknowledged the defect in Issue \#6373; while Pull Request \#6374 was closed, the maintainer unified sample allocation using \texttt{RL\_CALLOC} in commit 7639166~\cite{raylib_issue6373}.

\paragraph{Defect D1 (\drogon): Process Abort in \texttt{noexcept}.}
As shown in Figure~\ref{fig:motivating}, file saving operations in \drogon's \texttt{HttpFileImpl::save()} and \texttt{saveAs()} were marked \texttt{noexcept}.
However, the code invoked throwing overloads of standard filesystem functions such as \texttt{std::filesystem::exists}.
When an HTTP request uploaded a multipart file whose name exceeded 255 bytes, \texttt{std::filesystem::exists} threw \texttt{std::filesystem\_error} with \texttt{ENAMETOOLONG}.
Because the calling method was \texttt{noexcept}, the C++ runtime called \texttt{std::terminate()}, killing the entire server process.
The maintainers accepted our diagnosis and merged Pull Request \#2662, switching to non-throwing \texttt{std::error\_code} overloads~\cite{drogon_pr2662}.

\paragraph{Defect D2 (\drogon): Out-of-Bounds Iterator Dereference in Path Traversal Check.}
In \texttt{HttpFileImpl::save()}, path traversal checks compare the canonicalized upload path components against the target directory using a three-iterator loop.
When processing filenames containing directory traversal prefixes such as \texttt{../../c}, the weakly canonicalized path contained fewer path components than the upload root.
The comparison loop checked \texttt{*it1 != *it2} while advancing \texttt{++it1} without checking whether \texttt{it1} had reached its end boundary.
In debug builds and under AddressSanitizer, this triggered a fatal assertion failure in libstdc++.
We submitted an initial report in Pull Request \#2656, after which the maintainer committed the corresponding path component bounds check~\cite{drogon_pr2656}.

\paragraph{Additional Unfiled Medium- and High-Severity Defects.}
In addition to the six upstream resolved defects, \tool discovered ten medium- and high-severity defects that were verified with standalone crashing reproducers.
In \raylib, we uncovered two critical heap buffer overflows: an arbitrary write past allocated bone transformation matrices in \texttt{UpdateModelAnimationBones()} (CWE-122) and an audio stream buffer overflow of up to 76.8~KB in \texttt{UpdateAudioStreamInLockedState()} (CWE-122).
We also identified an out-of-bounds heap read in glTF accessor buffer views (CWE-125) and an unchecked mixer callback NULL dereference in \texttt{ReadAudioBufferFramesInInternalFormat()} (CWE-476).
Additional defects include an integer underflow in \texttt{ImageBlurGaussian()} (CWE-124) and a backward scan buffer underflow in \texttt{GetCodepointPrevious()} (CWE-124).
In \drogon, \tool detected premature header termination causing body desynchronization and request smuggling in \texttt{HttpRequestParser.cc} (CWE-444).
We also identified an out-of-range enum load under status codes above 599 triggering undefined behavior in \texttt{HttpResponseParser.cc} (CWE-843), a multipart boundary semantic desynchronization (RFC~2046), and header null-byte injection in \texttt{HttpUtils.cc} (CWE-113).

\begin{observation}
\label{obs:bugs}
\tool uncovered 16 security vulnerabilities across the three codebases: 6 acknowledged or merged upstream defects, and 10 additional medium- and high-severity defects.
Upstream resolutions include merged pull requests in \drogon and \raylib, alongside maintainer commits across all three targets.
\end{observation}

\subsection{\rqref{4}: Overhead and Cost}
\label{sec:rq4}

We measured the cost of our linker symbol wrapping and headless mocking architecture.
GNU linker symbol wrapping (\texttt{-Wl,--wrap}) introduces a single direct jump instruction per coverage edge event.
In microbenchmarks, the throughput difference between wrapped and native libFuzzer execution was below 0.8\%, within measurement noise.
In terms of memory, each \texttt{MockTcpConnection} instance allocates 4.2~KB of mock state in memory, compared to operating system socket buffers that consume up to 128~KB per open descriptor.
Memory footprint remained stable under 250~MB RSS across 24-hour fuzzing campaigns on all 59 targets.

\begin{observation}
\label{obs:overhead}
Linker-wrapped coverage interception imposes less than 1\% runtime overhead, and in-memory mock connections consume 4.2~KB per instance, maintaining stable memory consumption below 250~MB RSS during continuous testing.
\end{observation}
